\chapter{Criterio de ficha pública mínima de claim}

\section{Ficha pública mínima}
\begin{definition}
Se llama ficha pública mínima de claim al formato externo más sobrio mediante el cual un claim ya suficientemente estabilizado puede presentarse fuera del circuito interno de curación, conservando identidad, origen, alcance y reservas relevantes.
\end{definition}

\begin{observation}
La ficha pública mínima no sustituye la demo de origen, la tabla madre ni la auditoría interna. Su función es ofrecer una unidad externa clara, breve y disciplinada.
\end{observation}

\section{Condición mínima de fichabilidad pública}
\begin{definition}
Un claim puede adoptar ficha pública mínima cuando concurren simultáneamente:
\begin{enumerate}
  \item compatibilidad con publicación selectiva;
  \item trazabilidad suficientemente limpia hacia su origen;
  \item alcance sobrio y legible;
  \item posibilidad de distinguir soporte efectivo, interpretación y reserva;
  \item ausencia de ambigüedad inmediata sobre el estatuto editorial del claim.
\end{enumerate}
\end{definition}

\begin{observation}
La fichabilidad pública no exige demostrar todo el sistema en una sola pieza. Exige solo que la unidad externa no deforme la fuerza metodológica real del claim.
\end{observation}

\section{Campos mínimos de la ficha}
\begin{definition}
Toda ficha pública mínima debe poder declarar, al menos, los siguientes campos:
\begin{enumerate}
  \item identificador;
  \item formulación breve del claim;
  \item demo o fuente de origen;
  \item dependencia canónica visible;
  \item alcance declarado;
  \item condición mínima de validez;
  \item condición mínima de fallo;
  \item estado editorial;
  \item nota breve de reserva, si corresponde.
\end{enumerate}
\end{definition}

\begin{observation}
Estos campos no agotan toda presentación pública posible, pero bastan para fijar una primera unidad externa sobria y trazable.
\end{observation}

\section{Ficha y tabla madre}
\begin{definition}
Toda ficha pública mínima debe mantener vínculo legible con la tabla madre de demos y claims, de forma que su posición externa no borre su genealogía interna.
\end{definition}

\begin{observation}
La ficha pública no debe aparecer como fragmento autónomo sin historia metodológica.
\end{observation}

\section{Condición mínima de no fichabilidad}
\begin{definition}
Un claim no debe recibir ficha pública mínima cuando todavía depende de reservas internas no explicitadas, su formulación breve induce sobrelectura o su compresión externa destruye información crítica para su interpretación sobria.
\end{definition}

\begin{observation}
La no fichabilidad pública no degrada el claim. Puede significar solo que aún requiere anexo, reformulación o una vía externa distinta.
\end{observation}

\section{Cláusula de alcance}
\begin{observation}
El criterio de ficha pública mínima de claim fija la unidad externa más sobria del Tomo III. No constituye todavía una política completa de comunicación, dossier externo o publicación integral del corpus.
\end{observation}

